\documentclass[../main.tex]{subfiles}
\begin{document}
\section{Cifrado Asimétrico y Firma Digital}

Las clases anteriores resolvían confidencialidad (criptosistemas) e integridad
(MACs/hash) desde el modelo \textbf{simétrico}: emisor y receptor comparten una
\emph{misma} clave. Esta unidad revisita los mismos problemas desde el modelo
\textbf{asimétrico}, donde hay \emph{dos} claves distintas. Esto resuelve un
problema que antes era imposible (el intercambio de claves) y da origen a tres
construcciones nuevas: \textbf{intercambio de claves} (Diffie-Hellman),
\textbf{criptosistemas asimétricos} (RSA, ElGamal) y \textbf{firmas digitales}.

\prof{Hoy vamos a revisitar la criptografía, pero desde otro modelo mental que es el modelo asimétrico.}

\subsection{El problema de distribución de claves}

La seguridad de toda función criptográfica depende de que un atacante \textbf{no
tenga acceso a la clave}. Consecuencia práctica: las claves \textbf{no} pueden
transmitirse por un canal inseguro (la información cifrada sí, las claves no).
\begin{itemize}[leftmargin=1.4em]
  \item \textbf{Dos puntos:} se usa un canal alternativo seguro (encuentro
  físico, correo caro, o seguridad por capas con un canal independiente).
  \item \textbf{Múltiples puntos:} para $n$ entidades que se comunican de a pares
  hace falta una clave por combinación, del orden de $\binom{n}{2}=O(n^2)$ claves;
  cada parte administra $n-1$ claves. No escala (ej.: $100\,000$ empleados $\Rightarrow$
  $100\,000^2$ claves, cada uno conoce $999\,999$).
  \item \textbf{Punto único de confianza (TTP):} un tercero de confianza
  \emph{Trusted Third Party}, también llamado \textbf{KDC} (Key Distribution
  Center). Comparte una clave con cada entidad ($k_a,k_b,k_c,\dots$); si A quiere
  hablar con C, el KDC genera una clave de sesión $k_s$ y se la manda a cada uno
  cifrada con $k_a$ y $k_c$. Permite $n$ claves para $n$ entidades, pero es un
  \textbf{único punto de falla}. Base de Kerberos y Active Directory.
\end{itemize}

\subsection{La idea de la criptografía asimétrica}

Diffie y Hellman (1976, ``We stand today on the brink of a revolution in
cryptography'') parten de una analogía: hay problemas con naturaleza
\textbf{asimétrica}, fáciles en un sentido y difíciles en otro (cerrar un candado
es fácil, abrirlo sin llave es difícil; el problema \textbf{SAT} es difícil de
resolver pero fácil de verificar).

La idea: usar \textbf{dos claves independientes} (conocer una no permite derivar la
otra), una para cifrar y otra para descifrar. Si un atacante tiene la clave de
cifrado, igual no puede recuperar el mensaje. Entonces esa clave de cifrado se
puede \textbf{publicar a propósito} $\Rightarrow$ criptografía de \emph{clave
pública}. Esto resuelve elegantemente la distribución de claves: B publica su
clave de cifrado por un canal inseguro, A la usa para cifrar, y solo B (que tiene
la clave de descifrado) puede recuperar el mensaje.

\begin{destacado}
Propiedades de las claves (memorizar):
\begin{itemize}[leftmargin=1.4em,nosep]
  \item La clave \textbf{privada/secreta} se mantiene \textbf{secreta}, fuera del
  alcance de atacantes.
  \item La clave \textbf{pública} se \textbf{distribuye} (incluso al atacante),
  pero de forma \textbf{autenticada} (saber que es de quien dice ser).
  \item El uso de las claves \textbf{no es intercambiable}. En un criptosistema se
  cifra solo con la pública y se descifra solo con la privada. Usarlas al revés
  rompe la confidencialidad.
\end{itemize}
\end{destacado}

\begin{center}
\begin{tikzpicture}[
  font=\footnotesize,
  party/.style={draw, rounded corners, minimum width=2.1cm, minimum height=0.95cm, align=center, thick},
  arrow/.style={-{Stealth}, thick},
  key/.style={align=center, font=\scriptsize}
]
  % Simétrico
  \node[party] (sa) at (0,2.4) {Emisor};
  \node[party] (sb) at (5.2,2.4) {Receptor};
  \draw[arrow] (sa) -- node[above]{$c=\mathrm{Enc}_k(m)$} (sb);
  \node[key, below=2pt of sa] {clave $k$};
  \node[key, below=2pt of sb] {misma clave $k$};
  \node[font=\scriptsize\bfseries, above=10pt of sa.north east, xshift=1.0cm] {Simétrico (misma clave)};

  % Asimétrico
  \node[party] (aa) at (0,0) {Emisor};
  \node[party] (ab) at (5.2,0) {Receptor};
  \draw[arrow] (aa) -- node[above]{$c=\mathrm{Enc}_{pk_B}(m)$} (ab);
  \node[key, below=2pt of aa, text width=2.6cm] {usa la \textbf{pública} $pk_B$ del receptor};
  \node[key, below=2pt of ab, text width=2.6cm] {descifra con su \textbf{privada} $sk_B$};
  \node[font=\scriptsize\bfseries, above=10pt of aa.north east, xshift=1.0cm] {Asimétrico (dos claves)};
\end{tikzpicture}
\end{center}
\diagcaption{Simétrico: emisor y receptor comparten la misma clave $k$. Asimétrico:
el emisor cifra con la clave \emph{pública} del receptor y solo el receptor descifra
con su clave \emph{privada}.}

\paragraph{Problemas difíciles que la sustentan.} Toda esta criptografía se apoya
en problemas matemáticos sin solución eficiente conocida:
\begin{itemize}[leftmargin=1.4em,nosep]
  \item \textbf{Factorización} de $n=p\cdot q$ (sustenta RSA).
  \item \textbf{Logaritmo discreto} (DLP) en grupos finitos (sustenta DH y ElGamal).
\end{itemize}
``No eficiente'' significa que las mejores soluciones son exponenciales o
sub-exponenciales (equivalente a fuerza bruta).

\begin{ojo}
La computación cuántica afecta a \textbf{toda} esta parte asimétrica: el
\textbf{algoritmo de Shor} resolvería factorización y log discreto eficientemente
(pertenecen a la misma categoría de problemas). La criptografía simétrica clásica
no se ve afectada de la misma forma. Existen algoritmos post-cuánticos de
intercambio de claves para reemplazarlos.
\end{ojo}

\subsection{Repaso de álgebra y aritmética}

\begin{itemize}[leftmargin=1.4em]
  \item \textbf{Grupo} $(G,+)$: clausura, asociatividad, neutro, inverso.
  \textbf{Abeliano} si además es conmutativo. \textbf{Anillo} $(G,+,*)$: $(G,+)$
  abeliano $+$ clausura/asociatividad de $*$ $+$ distributividad. \textbf{Cuerpo/campo}:
  anillo con $*$ conmutativa, neutro e inverso de $*$.
  \item \textbf{Grupo cíclico} $G=(\{g^n\mid n\in\mathbb{Z}\},+)$. Grupo canónico
  $\mathbb{Z}_n=(\{0,1,\dots,n-1\},+)$. Un \textbf{generador} $g$ es coprimo con
  $n$; hay $\Phi(n)$ generadores. $g$ es \textbf{primitivo} si $\mathrm{ord}(g)=n$
  (recorre todos los elementos).
  \item \textbf{Aritmética modular:} $a \equiv b \bmod n \Leftrightarrow a-b=k\cdot n$.
  Si $n$ es primo $\Rightarrow$ campo finito $\mathbb{Z}_p$; si no, anillo (y
  $\mathbb{Z}_n^*$ con los coprimos a $n$ forma grupo multiplicativo).
  \item \textbf{Euler / Fermat:} si $(k:n)=1$ entonces $\exists\, k^{-1}$ tal que
  $k\cdot k^{-1}\equiv 1 \bmod n$, y
  \[ a^{\Phi(n)} \equiv 1 \bmod n, \]
  donde $\Phi(n)=\lvert\mathbb{Z}_n^*\rvert$. Caso primo (Fermat):
  $a^{p-1}\equiv 1 \bmod p$.
  \item \textbf{Cálculo de $\Phi$:} $\Phi(n\cdot m)=\Phi(n)\Phi(m)$ si $(n:m)=1$;
  $\Phi(p^k)=p^{k-1}(p-1)$ con $p$ primo. En RSA: $\Phi(p\cdot q)=(p-1)(q-1)$.
\end{itemize}

\subsection{Intercambio de claves}

\begin{definicion}{Protocolo de intercambio de claves}
Protocolo $\Pi(n)$ ejecutado por dos partes, sin entrada salvo el parámetro de
seguridad. Su salida es: una \textbf{transcripción} (mensajes intercambiados, vistos
por el atacante), una clave $k_a$ que deriva una parte y una clave $k_b$ que deriva
la otra, cada una con la información de su propio extremo. \textbf{Condición
fundamental:} $k_a=k_b$.
\end{definicion}

Tiene que ser la misma clave porque son dos entidades que no comparten nada de
antemano y luego necesitan una clave común para usar criptografía simétrica.

\paragraph{Seguridad (ataques pasivos): Key-Exchange experiment $KE_{A,\Pi}$.}
\begin{enumerate}[leftmargin=1.6em,nosep]
  \item Se ejecuta $\Pi$, sea $k=k_a=k_b$.
  \item Se genera $b\leftarrow\{0,1\}$.
  \item Si $b=0 \Rightarrow k'\leftarrow\{0,1\}^n$ (aleatoria); si $b=1 \Rightarrow k'=k$.
  \item $A$ recibe $Trans$ y $k'$, y emite $b'\in\{0,1\}$.
\end{enumerate}
$KE_{A,\Pi}=1$ si $b=b'$. Si $\Pr[KE_{A,\Pi}=1]<\tfrac12+\varepsilon$, $\Pi$ es
seguro. No se le pide al atacante \emph{adivinar} la clave, sino distinguir la
clave real de una aleatoria; así garantizamos que no obtiene \emph{ninguna}
información parcial. Adivinar al azar o sesgar da siempre $\approx 0.5$; si llega
a $\sim0.6$ es que encontró algo y el protocolo falla.

\subsubsection{Protocolo Diffie-Hellman}

\begin{enumerate}[leftmargin=1.6em]
  \item A define $(G,q,g)$: $G$ grupo, $q$ su tamaño, $g$ un generador. (Info
  \textbf{pública}, parte del transcript; un atacante la conoce.)
  \item A elige $x\leftarrow\mathbb{Z}_q$ y calcula $h_1=g^x$.
  \item A envía a B: $(G,q,g,h_1)$.
  \item B elige $y\leftarrow\mathbb{Z}_q$ y calcula $h_2=g^y$.
  \item B envía a A: $(h_2)$.
  \item A calcula $k_a=h_2^{\,x}=(g^y)^x$.
  \item B calcula $k_b=h_1^{\,y}=(g^x)^y$.
\end{enumerate}
La clave común (de sesión) es
\[ k_s=(g^x)^y \bmod p=(g^y)^x \bmod p=g^{xy}\bmod p, \]
e $k_a=k_b$ porque la exponenciación sucesiva es conmutativa. Lo que viaja son
$g^x$ y $g^y$; \textbf{$x$ nunca sale de A} y \textbf{$y$ nunca sale de B}.

\begin{center}
\begin{tikzpicture}[
  font=\footnotesize,
  arrow/.style={-{Stealth}, thick},
  box/.style={align=center}
]
  % columnas
  \node[box, font=\bfseries] (ah) at (0,3.4) {Alice};
  \node[box, font=\bfseries] (bh) at (7,3.4) {Bob};
  \draw[thick] (0,3.1) -- (0,-1.6);
  \draw[thick] (7,3.1) -- (7,-1.6);

  \node[box, anchor=north] (ax) at (0,3.0) {elige $x\leftarrow\mathbb{Z}_q$\\ calcula $g^x$};
  \node[box, anchor=north] (by) at (7,3.0) {elige $y\leftarrow\mathbb{Z}_q$\\ calcula $g^y$};

  \draw[arrow] (0,1.9) -- node[above]{$(G,q,g,\ g^x)$} (7,1.9);
  \draw[arrow] (7,1.0) -- node[above]{$g^y$} (0,1.0);

  \node[box, anchor=north] (ak) at (0,0.4) {$k_a=(g^y)^x$};
  \node[box, anchor=north] (bk) at (7,0.4) {$k_b=(g^x)^y$};

  \node[box, font=\bfseries] at (3.5,-0.9) {Clave común: $k_s=g^{xy}\bmod p$};

  % espía
  \node[draw, rounded corners, dashed, align=center, font=\scriptsize] (eve) at (3.5,2.0) {Espía ve\\ $g^x,\,g^y$\\ \textbf{no} obtiene $x,y,g^{xy}$ (DLP)};
\end{tikzpicture}
\end{center}
\diagcaption{Diffie-Hellman: Alice y Bob intercambian $g^x$ y $g^y$ y cada uno
deriva la misma clave $g^{xy}$. Un espía pasivo ve $g^x$ y $g^y$ pero no puede
recuperar $x$, $y$ ni $g^{xy}$ (problema del logaritmo discreto).}

\begin{examen}
RSA pregunta por DH casi siempre así:
\begin{itemize}[leftmargin=1.4em,nosep]
  \item \textbf{Detallar el protocolo:} $g$, $q$, $x\to g^x$, $y\to g^y$, clave
  común $g^{xy}$.
  \item \textbf{V/F ``DH permite ENVIAR una clave a Bob'' $\to$ FALSO.} DH
  \textbf{acuerda/deriva} una clave compartida \emph{en línea}; \textbf{no
  transporta} una clave preexistente. (2018/2023.) No confundir acuerdo de clave
  con transporte de clave.
  \item \textbf{Seguridad = problema del logaritmo discreto (DLP).} Dado $g^x$ y
  $g^y$ no se puede obtener $x$ ni $y$ (condición \emph{necesaria}). Formalmente
  requiere además la \textbf{conjetura de decisión DH}: dados $g,g^x,g^y$, un
  adversario no puede distinguir $g^{xy}$ de un valor aleatorio.
  \item \textbf{Valores válidos de $q$:} se trabaja con grupos multiplicativos, y
  para evitar líos con factores de primalidad se elige $q$ \textbf{primo}, así
  todos los elementos son coprimos. (Recup 2025, 2025-1.)
\end{itemize}
\end{examen}

\paragraph{Por qué es difícil.} Calcular logaritmos en campos \emph{infinitos} es
eficiente (logaritmo natural, series de Taylor). En campos \emph{finitos} (módulo
algo) el ``efecto de envolvente'' lo rompe: $3^2 \bmod 7 = 2$ (no 9), y recuperar
$x$ de $g^x \bmod p$ es el \textbf{problema del logaritmo discreto}, sin solución
eficiente (problema abierto de más de un siglo).

\begin{ojo}
\textbf{DH puro no tiene autenticación.} La versión original requiere un canal de
transmisión \textbf{autenticado}. Solo es seguro frente a atacantes
\textbf{pasivos}. Un atacante \textbf{activo} que modifica mensajes en tránsito
hace un ataque \textbf{Man-in-the-Middle (MiTM)} y obtiene las claves. Por eso en
la práctica DH se complementa con \textbf{MACs o firmas digitales}.
\end{ojo}

\begin{center}
\begin{tikzpicture}[
  font=\footnotesize,
  arrow/.style={-{Stealth}, thick},
  party/.style={draw, rounded corners, minimum width=1.6cm, minimum height=0.8cm, align=center, thick}
]
  \node[party] (a) at (0,0) {Alice\\ \scriptsize($x$)};
  \node[party, fill=red!10] (m) at (4,0) {Atacante\\ \scriptsize($m_1,m_2$)};
  \node[party] (b) at (8,0) {Bob\\ \scriptsize($y$)};

  \draw[arrow] (a) -- node[above]{$g^x$} (m);
  \draw[arrow] (m) -- node[above]{$g^{m_2}$} (b);
  \draw[arrow] (b) to[bend left=22] node[below]{$g^y$} (m);
  \draw[arrow] (m) to[bend left=22] node[below]{$g^{m_1}$} (a);

  \node[align=center, font=\scriptsize, below=1.2cm of a] {clave con atacante:\\ $g^{x m_1}$};
  \node[align=center, font=\scriptsize, below=1.2cm of b] {clave con atacante:\\ $g^{y m_2}$};
\end{tikzpicture}
\end{center}
\diagcaption{MiTM sobre DH puro (sin autenticación): el atacante intercepta y
negocia una clave distinta con cada parte ($g^{x m_1}$ con Alice, $g^{y m_2}$ con
Bob), descifrando y reenviando todo el tráfico.}

\subsection{Criptosistema asimétrico}

\begin{definicion}{Criptosistema asimétrico}
Terna de algoritmos:
\begin{align*}
\mathrm{Gen}&: ()\to PK\times SK &&\text{(genera \textbf{par} de claves pública/privada)}\\
\mathrm{Enc}&: PK\times P\to C &&\text{(cifra con la pública)}\\
\mathrm{Dec}&: SK\times C\to P &&\text{(descifra con la privada)}
\end{align*}
\textbf{Propiedad básica:} para todo $m$ y $(pk,sk)$ válidos, $d_{sk}(e_{pk}(m))=m$.
\end{definicion}

\paragraph{Prueba de indistinguibilidad ($Eav_{A,\Pi}$).} Igual que en el mundo
simétrico, pero \textbf{se le da $pk$ al adversario}: (1) $(pk,sk)\leftarrow K$;
(2) $A$ obtiene $pk$ y genera $(m_0,m_1)$; (3) $b\leftarrow\{0,1\}$; (4) $A$ recibe
$c=e_{pk}(m_b)$; (5) $A$ emite $b'$. Gana si $b=b'$; es indistinguible si
$\Pr[Eav_{A,\Pi}=1]<\tfrac12+\varepsilon$.

\begin{destacado}
Como $A$ conoce $pk$, puede cifrar cualquier mensaje (tiene acceso a la función de
cifrado, que es justo el oráculo del test CPA). Por eso, en el mundo asimétrico:
\[ \text{indistinguible (pasivo)} \;\Rightarrow\; \text{CPA-Secure} \]
y son pruebas equivalentes. Y como CPA-Secure \textbf{requiere cifrado no
determinístico}, \textbf{no existen criptosistemas asimétricos seguros que sean
determinísticos}: hay que incluir no-determinismo desde el principio.
\end{destacado}

\subsubsection{RSA (Rivest, Shamir, Adleman)}

Primer criptosistema asimétrico famoso. La parte compleja es la \textbf{generación
de claves}; cifrar y descifrar son simples.

\paragraph{``Textbook'' RSA (determinístico, solo conceptual).}
\begin{align*}
&\text{Elegir } p,q \text{ primos},\quad n=p\cdot q.\\
&e\leftarrow(0,\Phi(n)) \;\big|\; \mathrm{mcd}(e,\Phi(n))=1.\\
&\text{Calcular } d \;\big|\; e\cdot d\equiv 1 \bmod \Phi(n) \quad\text{(inverso multiplicativo, Euclides extendido)}.\\
&pk=(n,e),\qquad sk=(n,d).\\
&\mathrm{Enc}_{pk}(m)\equiv m^e \bmod n,\qquad \mathrm{Dec}_{sk}(c)\equiv c^d \bmod n.
\end{align*}
El cálculo $m^e \bmod n$ se hace con \textbf{exponenciación rápida}, reduciendo
módulo $n$ en cada paso, para no manejar números gigantes.

\begin{center}
\begin{tikzpicture}[
  font=\footnotesize,
  arrow/.style={-{Stealth}, thick},
  rsabox/.style={draw, rounded corners, align=center, minimum height=0.85cm, thick}
]
  % Generación
  \node[rsabox] (pq) at (0,0) {$p,q$ primos};
  \node[rsabox] (n) at (2.7,0) {$N=p\cdot q$\\ $\varphi(N)=(p{-}1)(q{-}1)$};
  \node[rsabox] (e) at (6.0,0) {$e$ con\\ $\mathrm{mcd}(e,\varphi(N))=1$};
  \node[rsabox] (d) at (9.3,0) {$d=e^{-1}$\\ $\bmod\,\varphi(N)$};
  \draw[arrow] (pq) -- (n);
  \draw[arrow] (n) -- (e);
  \draw[arrow] (e) -- (d);
  \node[align=center, font=\scriptsize, below=4pt of n] {$pk=(N,e)$};
  \node[align=center, font=\scriptsize, below=4pt of d] {$sk=(N,d)$};

  % Cifrado / descifrado
  \node[rsabox, fill=blue!6] (m1) at (1.0,-2.4) {$m$};
  \node[rsabox, fill=blue!6] (c) at (5.0,-2.4) {$c=m^e\bmod N$};
  \node[rsabox, fill=blue!6] (m2) at (9.3,-2.4) {$m=c^d\bmod N$};
  \draw[arrow] (m1) -- node[above, font=\scriptsize]{cifra con $pk$} (c);
  \draw[arrow] (c) -- node[above, font=\scriptsize]{descifra con $sk$} (m2);
\end{tikzpicture}
\end{center}
\diagcaption{RSA: arriba la generación de claves ($p,q\to N,\varphi(N),e,d$); abajo
cifrado $c=m^e\bmod N$ con la pública y descifrado $m=c^d\bmod N$ con la privada.}

\begin{examen}
\textbf{RSA es tema central de parcial.} Hay que saber de memoria:
\begin{itemize}[leftmargin=1.4em,nosep]
  \item Generación: $n=pq$, $\Phi(n)=(p-1)(q-1)$, elegir $e$ coprimo con $\Phi(n)$,
  $d=e^{-1}\bmod\Phi(n)$.
  \item Cifrado: $c=m^e \bmod n$. \qquad Descifrado: $m=c^d \bmod n$.
\end{itemize}
\end{examen}

\paragraph{Por qué funciona (Dec inversa de Enc).}
\begin{align*}
\mathrm{Dec}(\mathrm{Enc}(m)) &= (m^e)^d \bmod n = m^{ed}\bmod n.\\
\text{Como } e\cdot d\equiv 1 \bmod \Phi(n):\quad & e\cdot d = c\cdot\Phi(n)+1 \text{ (para alguna constante } c).\\
m^{ed} &= m^{c\Phi(n)+1} = \big(m^{\Phi(n)}\big)^{c}\cdot m \bmod n.\\
\text{Por Euler } m^{\Phi(n)}\equiv 1:\quad & = 1^c\cdot m = m \bmod n.
\end{align*}

\paragraph{Problemas del textbook RSA.}
\begin{itemize}[leftmargin=1.4em,nosep]
  \item \textbf{Es determinístico} $\Rightarrow$ no pasa la prueba de
  indistinguibilidad.
  \item \textbf{$e$ y $m$ chicos:} si $m^e<n$ nunca se reduce, y el log discreto
  coincide con el log natural (calculable). Históricamente se usaba $e=3$.
  \item \textbf{Módulos repetidos:} si dos pares de claves comparten el mismo $n$,
  se puede recuperar $n$ y luego la clave privada.
  \item Los primos $p,q$ deben ser ``primos fuertes'', no cualesquiera.
\end{itemize}

\paragraph{Ejemplo numérico (parámetros chicos).}
\begin{align*}
&p=2357,\ q=2551,\ n=p\cdot q=6\,012\,707,\quad \Phi(n)=(p-1)(q-1)=6\,007\,800.\\
&e=3\,674\,911 \text{ (al azar)},\quad d=422\,191 \text{ (Euclides ext.)}.\\
&\mathrm{Enc}(m=5\,234\,673)=5\,234\,673^{\,3\,674\,911}\bmod 6\,012\,707=3\,650\,502.\\
&\mathrm{Dec}(3\,650\,502)=3\,650\,502^{\,422\,191}\bmod 6\,012\,707=5\,234\,673.
\end{align*}

\paragraph{PKCS\,\#1 v1.5: RSA con padding aleatorio.} Estándar de RSA Labs que
hace RSA \textbf{no determinístico}. Sea $k$ la longitud de $n$ en bytes; solo se
cifran mensajes de hasta $k-11$ bytes (sea $D$ la longitud de $m$):
\[ m' = \texttt{00000000}\ \|\ \texttt{00000010}\ \|\ r\ \|\ \texttt{00000000}\ \|\ m, \]
donde $r$ son $k-D-3$ bytes \textbf{aleatorios distintos de 0} (para sacar el
padding al descifrar buscando el byte 00). Esto da $\ge 8$ bytes aleatorios
($\ge 2^{64}$ variantes) $\Rightarrow$ misma clave/mensaje produce salidas
distintas. Al descifrar se recupera $m'$ y se quita el padding.

\begin{examen}
\textbf{RSA con padding aleatorio} (PKCS\,\#1 v1.5) se cree que es
\textbf{CPA-Secure}, pero \textbf{no es CCA-Secure} (se encontraron ataques de
texto cifrado elegido; es manipulable, no da integridad ni es cifrado
autenticado). (2025-2.) Es decir: el padding aleatorio lo hace seguro frente a
CPA, no frente a CCA.
\end{examen}

\paragraph{Tamaño de claves y uso práctico.} El nivel de seguridad es relativo al
tamaño de los conjuntos (en RSA, $n=pq$). Como factorizar/log discreto son
sub-exponenciales (más eficientes que fuerza bruta), hacen falta claves grandes:
campos numéricos $n\ge 1024$ bits (hoy se recomiendan 1536 o 2048). RSA es 1-2
órdenes de magnitud más lento que AES y solo cifra mensajes acotados, por eso se
suele usar para \textbf{cifrar una clave simétrica} y con ella el mensaje largo
(lo mejor de los dos mundos).

\subsubsection{ElGamal (basado en DH)}

ElGamal es como un Diffie-Hellman ``desacoplado en el tiempo'' convertido en
criptosistema.
\begin{align*}
&\textbf{Gen: } \text{Seleccionar } G,q,g \text{ (campo $G$ de tamaño $q$); } x\leftarrow\mathbb{Z}_q,\ h=g^x.\\
&pk=(G,q,g,h),\qquad sk=(G,q,g,x).\\
&\mathrm{Enc}_{pk}(m): y\leftarrow\mathbb{Z}_q,\quad c=(c_1,c_2)=(g^y,\ h^y\cdot m).\\
&\mathrm{Dec}_{sk}(c)=c_2/c_1^{\,x}.
\end{align*}
$c_1=g^y$ es ``la segunda parte de DH''; $h^y=(g^x)^y=g^{xy}$ es la clave de sesión
que enmascara a $m$ (como un one-time pad). Al descifrar, $c_1^x=g^{yx}$ es esa
misma clave, y dividir la cancela.

\begin{itemize}[leftmargin=1.4em,nosep]
  \item \textbf{Resultado teórico:} si la prueba de decisión DH es difícil en $G$,
  ElGamal es \textbf{CPA-Secure}.
  \item \textbf{Diferencias con RSA:} ElGamal es \textbf{probabilístico} por
  naturaleza (elige $y$ al azar al cifrar); permite reutilizar $(G,q,g)$; \textbf{no
  está limitado a campos numéricos} (sirve con anillos de polinomios / campos de
  Galois con $2^n$ elementos, y con \textbf{curvas elípticas}, donde el problema de
  decisión DH es más complejo). En ElGamal el nivel de seguridad es relativo a
  $n=q$.
  \item \textbf{Curvas elípticas (ECC):} para seguridad equivalente a RSA-2048 ($\approx
  128$ bits) bastan $\approx 320$ bits (vs.\ 2048). Claves más chicas, operaciones
  más rápidas. Cuando un algoritmo de clave pública lleva ``EC'', suele ser ElGamal
  sobre curva elíptica.
\end{itemize}

\subsection{Firmas digitales}

Son el equivalente asimétrico de los MACs (su objetivo es la \textbf{integridad}),
pero con ventajas que justifican el nombre.

\begin{definicion}{Firma digital}
Terna de algoritmos:
\begin{align*}
\mathrm{Gen}&: ()\to PK\times SK\\
\mathrm{Sign}&: SK\times P\to S &&\text{(firma con la \textbf{privada})}\\
\mathrm{Vrfy}&: PK\times P\times S\to\{0,1\} &&\text{(verifica con la \textbf{pública})}
\end{align*}
\textbf{Propiedad básica:} $\mathrm{Vrfy}_{pk}(m,\mathrm{Sign}_{sk}(m))=1$.
\end{definicion}

\begin{destacado}
El rol de las claves está \textbf{al revés} que en el cifrado: se \textbf{firma con
la privada} y se \textbf{verifica con la pública}. Si firmáramos con la pública,
cualquiera podría firmar y se perdería la integridad.
\end{destacado}

\begin{center}
\begin{tikzpicture}[
  font=\footnotesize,
  arrow/.style={-{Stealth}, thick},
  party/.style={draw, rounded corners, minimum width=2.0cm, minimum height=0.9cm, align=center, thick}
]
  \node[party] (f) at (0,0) {Firmante};
  \node[party] (v) at (7.5,0) {Verificador};
  \draw[arrow] (f) -- node[above, align=center]{$(m,\ \sigma)$} (v);
  \node[align=center, font=\scriptsize, below=4pt of f, text width=3.0cm]
    {$\sigma=\mathrm{Sign}_{sk}\big(H(m)\big)$\\ usa su \textbf{privada} $sk$};
  \node[align=center, font=\scriptsize, below=4pt of v, text width=3.4cm]
    {$\mathrm{Vrfy}_{pk}(m,\sigma)\overset{?}{=}1$\\ usa la \textbf{pública} $pk$ del firmante};
\end{tikzpicture}
\end{center}
\diagcaption{Firma digital: el firmante firma $H(m)$ con su clave \emph{privada} y
cualquiera verifica con la \emph{pública}. Es el rol inverso al cifrado, donde se
cifra con la pública y se descifra con la privada.}

\paragraph{Ventajas sobre los MACs.}
\begin{itemize}[leftmargin=1.4em,nosep]
  \item \textbf{Públicamente verificables:} cualquiera con la clave pública verifica
  (un MAC necesita la clave secreta para verificar).
  \item \textbf{Transferibles:} se pueden reenviar a varios destinatarios o por una
  cadena y siguen siendo verificables en cada eslabón.
  \item \textbf{No repudio:} quien firma no puede negar haberlo hecho. Como la única
  forma de generar la firma es con la clave privada (y se asume que nunca se
  comparte), tiene valor legal análogo a la firma de puño y letra (en Argentina hay
  ley de firma digital; un organismo del Estado valida que el dueño de la clave
  privada es tal persona real).
\end{itemize}

\paragraph{Seguridad: $\mathrm{Sig\text{-}forge}_{A,\Pi}$ (falsificación).}
Análoga a Mac-forge: (1) $(sk,pk)\leftarrow K$; (2) $A$ obtiene
$f(x)=\mathrm{Sign}_{sk}(x)$ y \textbf{$pk$}; (3) $A$ hace las consultas de firma que
quiera (conjunto $Q$); (4) $A$ emite $(m,s)$ con $m\notin Q$.
$\mathrm{Sig\text{-}forge}_{A,\Pi}=1$ si $\mathrm{Vrfy}_{pk}(m,s)=1$ y $m\notin Q$.
Si ningún $A$ gana, la firma es \textbf{infalsificable}.

\subsubsection{RSA-Signature y por qué se firma el HASH}

Idéntico a RSA-Encryption pero \textbf{invirtiendo los roles de las claves}:
\[ \mathrm{Sign}_{sk}(m)\equiv m^d \bmod n, \qquad
   \mathrm{Vrfy}_{pk}(m,s):\ \overset{?}{m=s^e \bmod n}. \]

\begin{ojo}
\textbf{RSA-Signature básico (sin hash) es INSEGURO.} Ataques de falsificación:
\begin{itemize}[leftmargin=1.4em,nosep]
  \item Elegir $s\leftarrow S$ al azar, calcular $m=s^e\bmod n$, emitir $(m,s)$:
  \textbf{pasa la verificación}.
  \item \textbf{Maleabilidad:} con $s_1=f(m_1)$, $s_2=f(m_2)$, el par
  $(m_1\!\cdot m_2,\ s_1\!\cdot s_2)$ también verifica.
\end{itemize}
\end{ojo}

\paragraph{Hashed RSA (la firma que se usa).} Solución elegante: en vez de firmar
el mensaje, firmar su \textbf{hash} con una función de hash criptográfica (libre de
colisiones):
\[ \mathrm{Sign}_{sk}(m)\equiv H(m)^d \bmod n, \qquad
   \mathrm{Vrfy}_{pk}(m,s):\ \overset{?}{H(m)=s^e \bmod n}. \]

\begin{examen}
\textbf{Firma RSA en el parcial (2015):}
\begin{itemize}[leftmargin=1.4em,nosep]
  \item Firmar: $\sigma=H(m)^d \bmod n$.
  \item Verificar (Vrfy): recibir $m$ y $\sigma$, calcular $\sigma^e \bmod n$ y
  \textbf{comparar con $H(m)$}.
  \item \textbf{Por qué se firma el HASH y no el mensaje:} (1) la resistencia a
  preimagen del hash evita el primer ataque (no se puede partir de $s$ y construir
  un $m$ válido); (2) no poder predecir $H(m)$ evita el ataque de maleabilidad;
  (3) bonus: el hash es de tamaño fijo, así la firma tiene tamaño fijo
  independiente del mensaje (RSA básico solo firma mensajes chicos del tamaño del
  bloque).
\end{itemize}
Hashed RSA no tiene prueba de seguridad salvo asumiendo un modelo ideal de $H$.
\end{examen}

\paragraph{DSS (Digital Signature Standard).} Algoritmo de firma basado en ElGamal
(usa $H$ = SHA1/SHA2). Genera claves $(p,q,g,y)$ / $(p,q,g,x)$ con $q$ primo,
$p$ primo tal que $q\mid p-1$, $g$ generador de orden $q$, $x\leftarrow\mathbb{Z}_q$,
$y=g^x\bmod p$. Firma $(r,s)$ con $r=(g^k\bmod p)\bmod q$ y
$s=[H(m)+x\cdot r]\cdot k^{-1}\bmod q$. Hereda las ventajas de ElGamal sobre RSA:
existe sobre curvas elípticas (ECDSA) $\Rightarrow$ firmas más compactas con claves
más chicas.

\begin{examen}
\textbf{Firma digital para distribución de software / binarios} (2017): el
proveedor firma el binario con su clave privada; cualquiera verifica con la clave
pública que el binario es auténtico y no fue modificado. Aprovecha que la firma es
\emph{públicamente verificable} y da integridad/autenticación de origen sin
compartir secretos.
\end{examen}

\begin{ojo}
\textbf{Diferencia entre cifrar y firmar (no confundir):}
\begin{itemize}[leftmargin=1.4em,nosep]
  \item \textbf{Cifrar} $\to$ confidencialidad. Se cifra con la \textbf{pública}
  del destinatario, descifra con su privada.
  \item \textbf{Firmar} $\to$ integridad / autenticación / no repudio
  (\emph{no} da confidencialidad). Se firma con la \textbf{privada} del emisor,
  verifica cualquiera con la pública.
\end{itemize}
\end{ojo}

\subsection{PKI y certificados digitales}

El intercambio asimétrico resuelve la confidencialidad pero traslada el problema:
hay que distribuir la clave pública \textbf{de forma autenticada} (saber que la
pública de Bob es realmente de Bob, si no, un MiTM publica la suya). Esto lo
resuelve la \textbf{PKI} (Public Key Infrastructure) con \textbf{certificados} y
\textbf{Autoridades de Certificación (CA)}.

\begin{itemize}[leftmargin=1.4em]
  \item Un \textbf{certificado digital} asocia una identidad (titular) con su
  \textbf{clave pública}, y está \textbf{firmado por la CA} (con la clave privada
  de la CA).
  \item \textbf{Cadena de confianza:} se confía en la CA (cuya clave pública se
  conoce de antemano); la CA firma el certificado del titular. Pueden encadenarse
  (CA raíz $\to$ CA intermedia $\to$ titular).
\end{itemize}

\begin{center}
\begin{tikzpicture}[
  font=\footnotesize,
  arrow/.style={-{Stealth}, thick},
  cert/.style={draw, rounded corners, align=center, minimum width=2.4cm, minimum height=1.0cm, thick}
]
  \node[cert, fill=green!8] (root) at (0,0) {\textbf{CA raíz}\\ \scriptsize confianza\\ \scriptsize preinstalada};
  \node[cert] (inter) at (4.2,0) {\textbf{CA intermedia}\\ \scriptsize $pk$ + firma};
  \node[cert] (site) at (8.4,0) {\textbf{Cert. del sitio}\\ \scriptsize $pk$ del titular};

  \draw[arrow] (root) -- node[above, font=\scriptsize, align=center]{firma con\\ $sk_{\text{raíz}}$} (inter);
  \draw[arrow] (inter) -- node[above, font=\scriptsize, align=center]{firma con\\ $sk_{\text{interm}}$} (site);

  % validación hacia arriba
  \draw[arrow, dashed] (site.south) to[bend left=28] node[below, font=\scriptsize]{verifica firma con $pk_{\text{interm}}$} (inter.south);
  \draw[arrow, dashed] (inter.south) to[bend left=28] node[below, font=\scriptsize]{verifica firma con $pk_{\text{raíz}}$} (root.south);
\end{tikzpicture}
\end{center}
\diagcaption{Cadena de certificados: cada nivel firma el certificado del siguiente
con su clave privada. Validar = verificar las firmas hacia arriba hasta la CA raíz,
en la que se confía de antemano.}

\begin{examen}
\textbf{Certificados en el parcial:}
\begin{itemize}[leftmargin=1.4em,nosep]
  \item \textbf{Validar un certificado = verificar la FIRMA de la CA} (con la clave
  pública de la CA). (2018, 2017.)
  \item Un certificado contiene la clave \textbf{PÚBLICA del titular}, \textbf{NO}
  la clave privada de la CA. \textbf{V/F ``el certificado contiene la clave privada
  de la CA'' $\to$ FALSO.} (2025-1, 2025-2.)
  \item Propiedades clave pública/privada: la privada se mantiene secreta, la
  pública se distribuye de forma \textbf{autenticada} (justamente para eso sirve el
  certificado). (2015, 2025-1.)
\end{itemize}
\end{examen}

\begin{ojo}
\textbf{Errores típicos de esta unidad:}
\begin{itemize}[leftmargin=1.4em,nosep]
  \item Pensar que el certificado contiene la clave \emph{privada} (contiene la
  pública del titular, firmada por la CA).
  \item Creer que \textbf{MD5 es asimétrico}: MD5 es una \textbf{función de hash},
  no un esquema de clave pública.
  \item Confundir \textbf{DH (acuerdo de clave)} con \textbf{transporte de clave}:
  DH deriva una clave común, no envía una clave preexistente.
\end{itemize}
\end{ojo}

\end{document}
